\documentclass[11pt]{article}
\pagestyle{plain}

% \pgfplotsset{compat=1.18}
\usepackage[margin=1in]{geometry}
\usepackage{amsmath,amsfonts,amssymb,amsthm}
\usepackage[ruled,vlined,linesnumbered,commentsnumbered]{algorithm2e}
\SetKwInput{KwIn}{Input}
\SetKwInput{KwOut}{Output}
\SetKwInput{KwGlobal}{Global variable}
\SetKwInput{KwOracle}{Oracle}
\SetKwComment{Comment}{// }{}
\usepackage{graphicx}
\usepackage{enumerate}
\usepackage{bbm}
\usepackage{verbatim}
\usepackage{hyperref,color}
\usepackage[capitalize,nameinlink]{cleveref}
\usepackage[dvipsnames]{xcolor}
\hypersetup{
	colorlinks=true,
	pdfpagemode=UseNone,
	citecolor=OliveGreen,
	linkcolor=NavyBlue,
	urlcolor=Magenta,
	pdfstartview=FitW
}
\usepackage{appendix}
\usepackage{makecell}
\usepackage{tablefootnote}
\crefname{appsec}{Appendix}{Appendices}
\usepackage{tikz}
\usepackage{pgfplots}
\usepackage{xifthen}
\usepackage{subcaption}
\usepackage{hhline}

\theoremstyle{plain}
\newtheorem{theorem}{Theorem}[section]
\newtheorem{proposition}[theorem]{Proposition}
\newtheorem{lemma}[theorem]{Lemma}
\newtheorem{corollary}[theorem]{Corollary}
\newtheorem{conjecture}[theorem]{Conjecture}
\newtheorem{problem}[theorem]{Problem}
\newtheorem{claim}[theorem]{Claim}
\newtheorem{fact}[theorem]{Fact}

\theoremstyle{definition}
\newtheorem{definition}[theorem]{Definition}
\newtheorem{example}[theorem]{Example}
\newtheorem{observation}[theorem]{Observation}
\newtheorem{assumption}[theorem]{Assumption}
\newtheorem*{assumption*}{Assumption}
\newtheorem{condition}[theorem]{Condition}

\theoremstyle{remark}
\newtheorem{remark}[theorem]{Remark}

\crefname{lemma}{Lemma}{Lemmas}
\crefname{theorem}{Theorem}{Theorems}
\crefname{definition}{Definition}{Definitions}
\crefname{fact}{Fact}{Facts}
\crefname{claim}{Claim}{Claims}
\crefname{proposition}{Proposition}{Propositions}


\newcommand{\dif}{\,\mathrm{d}}
%\newcommand{\E}{\mathbb{E}}
%\newcommand{\Var}{\mathrm{Var}}
\newcommand{\Cov}{\mathrm{Cov}}
\newcommand{\MI}{\mathrm{I}}
%\newcommand{\Ent}{\mathrm{Ent}}
\newcommand{\ind}{\mathbbm{1}}
\newcommand{\allone}{\mathbf{1}}
\newcommand{\tr}{\mathrm{tr}}
\newcommand{\wt}{\mathrm{weight}}
\DeclareMathOperator*{\argmax}{arg\,max}
\newcommand{\norm}[1]{\left\lVert #1 \right\rVert}
\newcommand{\TV}[2]{\left\|#1 - #2\right\|_{\mathrm{TV}}}
\newcommand{\ceil}[1]{\left\lceil #1 \right\rceil}
\newcommand{\floor}[1]{\left\lfloor #1 \right\rfloor}
\newcommand{\ip}[2]{\left\langle #1 , #2 \right\rangle}
\newcommand{\grad}{\nabla}
\newcommand{\hessian}{\nabla^2}
\newcommand{\trans}{\intercal}
\newcommand{\diag}{\mathrm{diag}}
\newcommand{\poly}{\mathrm{poly}}
\newcommand{\dist}{\mathrm{dist}}
\newcommand{\sgn}{\mathrm{sgn}}
\newcommand{\fpras}{\mathsf{FPRAS}}
\newcommand{\fpaus}{\mathsf{FPAUS}}
\newcommand{\fptas}{\mathsf{FPTAS}}

\newcommand{\eps}{\varepsilon}

\newcommand{\N}{\mathbb{N}}
\newcommand{\Z}{\mathbb{Z}}
\newcommand{\Q}{\mathbb{Q}}
\newcommand{\R}{\mathbb{R}}
\newcommand{\C}{\mathbb{C}}

\renewcommand{\AA}{\mathcal{A}}
\newcommand{\BB}{\mathcal{B}}
\newcommand{\CC}{\mathcal{C}}
\newcommand{\DD}{\mathcal{D}}
\newcommand{\EE}{\mathcal{E}}
\newcommand{\FF}{\mathcal{F}}
\newcommand{\GG}{\mathcal{G}}
\newcommand{\HH}{\mathcal{H}}
\newcommand{\II}{\mathcal{I}}
\newcommand{\JJ}{\mathcal{J}}
\newcommand{\KK}{\mathcal{K}}
\newcommand{\LL}{\mathcal{L}}
\newcommand{\MM}{\mathcal{M}}
\newcommand{\NN}{\mathcal{N}}
\newcommand{\OO}{\mathcal{O}}
\newcommand{\PP}{\mathcal{P}}
\newcommand{\QQ}{\mathcal{Q}}
\newcommand{\RR}{\mathcal{R}}
\renewcommand{\SS}{\mathcal{S}}
\newcommand{\TT}{\mathcal{T}}
\newcommand{\UU}{\mathcal{U}}
\newcommand{\VV}{\mathcal{V}}
\newcommand{\WW}{\mathcal{W}}
\newcommand{\XX}{\mathcal{X}}
\newcommand{\YY}{\mathcal{Y}}
\newcommand{\ZZ}{\mathcal{Z}}

\newcommand{\KL}[2]{D_{\mathrm{KL}} \left( #1 \,\Vert\, #2 \right)}
\newcommand{\KLbig}[2]{D_{\mathrm{KL}} \left( #1 \,\big\Vert\, #2 \right)}
\newcommand{\KLBig}[2]{D_{\mathrm{KL}} \left( #1 \,\Big\Vert\, #2 \right)}

\newcommand{\SKL}[2]{D_{\mathrm{SKL}} \left( #1 \,\Vert\, #2 \right)}
\newcommand{\SKLbig}[2]{D_{\mathrm{SKL}} \left( #1 \,\big\Vert\, #2 \right)}
\newcommand{\SKLBig}[2]{D_{\mathrm{SKL}} \left( #1 \,\Big\Vert\, #2 \right)}

\newcommand{\qandq}{\quad\text{and}\quad}
\newcommand{\supp}{\mathrm{supp}}

\newcommand{\DKL}[2]{\-D_{\-{KL}}\left(#1 \parallel #2\right)}
\newcommand{\DTV}[2]{\-D_{\mathrm{TV}}\left({#1},{#2}\right)}
% \newcommand{\dist}{\mathrm{dist}}
\newcommand{\e}{\mathrm{e}}
\renewcommand{\epsilon}{\varepsilon}
\renewcommand{\emptyset}{\varnothing}
% \newcommand{\norm}[1]{\left\Vert#1\right\Vert}
\newcommand{\set}[1]{\left\{#1\right\}}
\newcommand{\tuple}[1]{\left(#1\right)} 
% \newcommand{\eps}{\varepsilon}
\newcommand{\inner}[2]{\left\langle #1,#2\right\rangle}
\newcommand{\tp}{\tuple}
\newcommand{\ol}{\overline}
\newcommand{\abs}[1]{\left\vert#1\right\vert}
\newcommand{\ctp}[1]{\left\lceil#1\right\rceil}
\newcommand{\ftp}[1]{\left\lfloor#1\right\rfloor}
\newcommand{\Ind}[1]{\-{Ind}_{#1}}
\newcommand{\todo}[1]{{\color{red} TODO: #1}}

\newcommand{\Cor}{\Psi^{\mathrm{sym}}}
\newcommand{\cor}{\Psi^{\-{cor}}}
\newcommand{\Inf}{\Psi}
% \newcommand{\diag}{\-{diag}}

\def\*#1{\boldsymbol{#1}} % Use \*A for \mathbf{A}
\def\+#1{\mathcal{#1}} % Use \+A for \mathcal{A}
\def\-#1{\mathrm{#1}} % Use \-A for \mathrm{A}
\def\=#1{\mathbb{#1}} % Use \^A for \mathbb{A}
\def\!#1{\mathfrak{#1}} % Use \!A for \mathfrak{A}

\newcommand*\diff{\mathop{}\!\mathrm{d}}

% probability
% \DeclareMathOperator*{\oPr}{\mathbf{Pr}}
\def\oPr{\mathop{\mathrm{Pr}}}
\renewcommand{\Pr}[2][]{ \ifthenelse{\isempty{#1}}
  {\oPr\left[#2\right]}
  {\oPr_{#1}\left[#2\right]} } % Use \Pr[a]{b} for \mathbf{Pr}_a[b], \Pr{b} for  \mathbf{Pr}[b]

% expectation: relies on the package "dsfont"
% \DeclareMathOperator*{\oE}{\mathbb{E}}
\def\oE{\mathop{\mathbb{E}}}
\newcommand{\E}[2][]{ \ifthenelse{\isempty{#1}}
  {\oE\left[#2\right]}
  {\oE_{#1}\left[#2\right]} }

% variance
%\DeclareMathOperator*{\oVar}{\mathbf{Var}}
\def\oVar{\mathrm{Var}}
\newcommand{\Var}[2][]{ \ifthenelse{\isempty{#1}}
  {\oVar\left[#2\right]}
  {\oVar_{#1}\left[#2\right]} }

% entropy
% \DeclareMathOperator*{\oEnt}{\mathbf{Ent}}
\def\oEnt{\mathrm{Ent}}
\newcommand{\Ent}[2][]{ \ifthenelse{\isempty{#1}}
  {\oEnt\left[#2\right]}
  {\oEnt_{#1}\left[#2\right]} }

\newcommand{\PhiEnt}[2][]{ \ifthenelse{\isempty{#1}}
  {\oEnt^\phi\left[#2\right]}
  {\oEnt^\phi_{#1}\left[#2\right]} }


\newenvironment{Exampleparagraph}[1][]
{
    \par 
    \noindent 
    \textbf{\color{blue} #1}
    \par 
    \vspace{0.5em} 
    \itshape 
}
{
    \par 
    \vspace{1em} 
}

%\input{local-macro}

\newcommand{\RED}[1]{\textcolor{red}{#1}}
\newcommand{\zc}[1]{\textcolor{red}{ZC: #1}}

\newcommand{\bern}[1]{\mathrm{Bern}\left(#1\right)}



\title{Note on \hyperlink{https://arxiv.org/pdf/2407.17864}{https://arxiv.org/pdf/2407.17864}}
\date{\today}

\pgfplotsset{compat=1.18} 
\begin{document}

\maketitle

\section{Gaussian isoperimetric inequality}
Let $\gamma_n$ be $n$-dimensional standard Gaussian measure.

\begin{lemma}[Bobkov inequality]
    For any function $f: \mathbb{R}^n \to [0,1]$, it holds
    \begin{align*}
        I\tp{\int f \dif \gamma_n} \le \int \sqrt{I(f)^2 + \abs{\nabla f}^2} \dif \gamma_n,
    \end{align*}
    where $I(\cdot) = \phi(\Phi^{-1}(\cdot))$ is the Gaussian isoperimetric profile.
\end{lemma}
\begin{proof}
    Let $P_t$ be the Ornstein–Uhlenbeck semigroup (corresponding to the Lagenvin diffusion).
    It suffices to show that $\int F_t \dif \gamma_n$ is monotone-decreasing in $t$, where 
    \begin{align*}
        F_t(x) := \sqrt{I(P_tf(x))^2 + \abs{\nabla P_tf(x)}^2} (=: \sqrt{q}). 
    \end{align*}
    We prove by showing that $(\partial_t - L) F_t \le 0$ for all $x$, where $L = \Delta - x \cdot \nabla$ is the generator.

    By chain rule:
    \begin{align*}
        \nabla \sqrt{q} = \frac{1}{2q^{1/2}} \cdot \nabla q \text{ and }  \Delta \sqrt{q} = \frac{1}{2q^{1/2}} \cdot \Delta q - \frac{1}{4q^{3/2}} \cdot \abs{\nabla q}^2.
    \end{align*}
    Therefore, it holds
    \begin{align*}
        L \sqrt{q} = \frac{1}{2q^{1/2}} \cdot Lq - \frac{1}{4q^{3/2}} \cdot \abs{\nabla q}^2.
    \end{align*}
    It remains to show that
    \begin{align*}
    \abs{\nabla q}^2 \le 2q \cdot (L - \partial_t)q. 
    \end{align*}
    Let $u := P_t f$. Again by chain rule and identity $(\partial_t - L)\Gamma(u) = -2\Gamma_2(u)$, we have
    \begin{align*}
        (L-\partial_t) q = (L-\partial_t) (I(u)^2 + \abs{\nabla u}^2) = 2 \Gamma_2(u) + 2 (I'(u)^2 - 1) \Gamma(u),
    \end{align*}
    where $\Gamma(f,g) = \frac{1}{2} \tp{L(fg) - fL(g)-gL(f)}$ and $\Gamma_2(f) = \frac{1}{2}\tp{f \Gamma(f) - 2\Gamma\tp{f,Lf}}$.
    Also, by a straightforward calculation, it holds
    \begin{align*}
        \Gamma(q) &= \abs{\nabla q}^2 = 4 \abs{(\nabla^2 u) \nabla u + I(u) I'(u) \cdot \nabla u }^2 \le 4 \tp{\abs{\nabla^2 u}_{F} \abs{\nabla u} + \abs{I(u)} \abs{I'(u) \nabla u}}^2\\
        &\le 4 q \cdot \tp{\abs{\nabla^2 u}^2_F + \abs{I'(u) \nabla u}^2} = 4q \cdot \tp{\Gamma_2(u) + (I'(u)^2 - 1)\Gamma(u)}.
    \end{align*}
    This concludes the proof.
\end{proof}

\begin{corollary}[Gaussian isoperimetric inequality]
    Let $A \subseteq \mathbb{R}^n$ be a measurable set with volume $V=\gamma_n(A)$. The perimeter of set $A$  defined by $P(A) := \lim_{\eps \to 0^+} \frac{\gamma_n(A+B(0,\eps)) - \gamma_n(A)}{\eps}$ satisfies
    \begin{align*}
        P(A) \ge I(V).
    \end{align*}
\end{corollary}

\section{Gaussian isoperimetric inequality for general probability measures}

Recall $\Gamma(f,g) = \frac{1}{2} \tp{L(fg) - fL(g)-gL(f)}$ and $\Gamma_2(f) = \frac{1}{2}\tp{L\Gamma(f) - 2\Gamma(f,Lf)}$.

We let $\Gamma(f) = \Gamma(f,f)$ and $\Gamma_2(f) = \Gamma_2(f,f)$ for simplicity.

We assume that $(P_t)_{t \ge 0}$ is a diffusion semigroup, i.e., for any $f$ and smooth function $\phi$, 
\begin{align*}
    L(\phi(f)) = \phi'(f) L f + \phi''(f) \Gamma(f).
\end{align*}

\begin{fact}
    It holds
    \begin{align*}
        \Gamma(\phi(f),g) = \phi'(f) \Gamma(f,g).
    \end{align*}
\end{fact}
\begin{proof}
    By definition, it holds
    \begin{align*}
        \Gamma(\phi(f),g) &= \frac{1}{2} \tp{L(\phi(f) g) - \phi(f) L g - g L \phi(f)}\\
        &= \frac{1}{2} \tp{\phi'(f) g L f + \phi''(f) g \Gamma(f) + \phi(f) L g - \phi(f) L g - g \phi'(f) L f - g \phi''(f) \Gamma(f)}\\
        &= \phi'(f) \Gamma(f,g).
    \end{align*}
\end{proof}
\begin{fact}
    Let $f_t = P_t f$. It holds
    \begin{align*}
        (\partial_t - L) \Gamma(f_t) = -2 \Gamma_2(f_t).
    \end{align*}
\end{fact}
\begin{proof}
    By definition, it holds
    \begin{align*}
        (\partial_t - L) \Gamma(f_t) &= \frac{1}{2} \tp{(\partial_t - L) L(f_t^2) - 2 (\partial_t - L)(f_t L f_t)}\\
        &= \frac{1}{2} \tp{L^2 (f_t^2) - 2L(f_t L f_t)} = -2 \Gamma_2(f_t).
    \end{align*}
\end{proof}

In the following, we assume $\Gamma_2(f) \ge \rho \Gamma(f)$ and $\Gamma(\Gamma(f)) \le 4 \Gamma(f) \tp{\Gamma_2(f) - \rho \Gamma(f)}$.
\begin{lemma}[Bobkov inequality for general probability measures]
    For any function $f: \mathbb{R}^n \to [0,1]$, it holds
    \begin{align*}
        I\tp{\int f \dif \mu} \le \int \sqrt{I(f)^2 + \frac{1}{\rho} \cdot \Gamma(f)} \dif \mu,
    \end{align*}
    where $\mu$ is a probability measure satisfying $\Gamma_2(f) \ge 0$ for all $f$ and $\rho > 0$.
\end{lemma}
\begin{proof}
    Similar to the first proof, it suffices to show that $(\partial_t - L) F_t \le 0$ for all $x$, where
    \begin{align*}
        F_t(x) := \sqrt{I(f_t(x))^2 + \frac{1}{\rho} \cdot \Gamma(f_t)(x)} (=: \sqrt{q_t}).
    \end{align*}
    Note that
    \begin{align*}
        \partial_t \sqrt{q_t} = \frac{1}{2 \sqrt{q_t}} \cdot \partial_t q_t \text{ and } L \sqrt{q_t} = \frac{1}{2\sqrt{q_t}}\tp{L q_t - \frac{1}{2q_t} \Gamma(q_t)}.
    \end{align*}
    Therefore, it suffices to show that
    \begin{align*}
        \Gamma(q_t) \le 2 q_t (L - \partial_t) q_t.
    \end{align*}
    Note that
    \begin{align*}
        (L - \partial_t) I^2(f_t) &= 2 (I^2)'(f_t) \cdot (L - \partial_t) f_t + 2 (I^2)''(f_t) \Gamma(f_t)\\
        &= 2 I(f_t) \cdot I'(f_t) \cdot (L - \partial_t) f_t + 2 (I'(f_t)^2 -1) \Gamma(f_t)\\
        &= 2(I'(f_t)^2 -1) \Gamma(f_t).
    \end{align*}
    Therefore,
    \begin{align*}
        (L - \partial_t) q_t &= (L - \partial_t) I^2(f_t) + \frac{1}{\rho} (L - \partial_t) \Gamma(f_t)\\
        &= 2(I'(f_t)^2 -1) \Gamma(f_t) + \frac{2}{\rho} \Gamma_2(f_t)\\
        &= 2I'(f_t)^2 \Gamma(f_t) + \frac{2}{\rho} R_t,
    \end{align*}
    where we let $R_t = \Gamma_2(f_t) - \rho \Gamma(f_t) \ge 0$.
    Also, by a straightforward calculation, it holds
    \begin{align*}
        \Gamma(q_t) &= \Gamma(I^2(f_t)) + \frac{1}{\rho^2} \Gamma(\Gamma(f_t)) + \frac{2}{\rho} \Gamma(I^2(f_t), \Gamma(f_t))\\
        &= (I^2)'(f_t)^2 \Gamma(f_t) + \frac{1}{\rho^2} \Gamma(\Gamma(f_t)) + \frac{2}{\rho} (I^2)'(f_t) \Gamma(f_t, \Gamma(f_t))\\
        &= 4 I^2(f_t) I'(f_t)^2 \Gamma(f_t) + \frac{1}{\rho^2} \Gamma(\Gamma(f_t)) + \frac{4}{\rho} I(f_t) I'(f_t) \Gamma(f_t, \Gamma(f_t))\\
        &\le 4 I^2(f_t) I'(f_t)^2 \Gamma(f_t) + \frac{4}{\rho} I(f_t) I'(f_t) \Gamma(f_t, \Gamma(f_t)) + \frac{4}{\rho^2} \Gamma(f_t) R_t.
    \end{align*}
    It suffices to show that
    \begin{align*} 
        (I'(f_t))^2 \Gamma(f_t)^2+I^2(f_t) R_t \ge I(f_t) I'(f_t) \Gamma(f_t,\Gamma(f_t))
    \end{align*}
    This holds by Cauchy's inequality $\Gamma(f,g) \le \sqrt{\Gamma(f) \Gamma(g)}$.
\end{proof}

\paragraph{Application}
    Let $\mu(x) \propto \exp\tp{-V(x)}$ be an $\rho$-strongly log-concave distribution. Let $(P_t)_{t \ge 0}$ be the Markov semigroup for the Langevin dynamics.
    The infinitesimal generator $L = \Delta - \nabla V \cdot \nabla$,
    \begin{align*}
        \Gamma(f) = \abs{\nabla f}^2 \text{ and } \Gamma_2(f) = \|\hessian f\|_F^2 + \inner{\nabla^2 V \cdot \nabla f}{\nabla f} \ge \rho \Gamma(f).
    \end{align*}
    Moreover,
    \begin{align*}
        \Gamma(\Gamma f) = \abs{\nabla \abs{\nabla f}^2}^2 = 4 \abs{\nabla^2 f \cdot \nabla f}^2 \le 4 \abs{\nabla^2 f}_F^2 \cdot \abs{\nabla f}^2 \le 4 \Gamma(f) \tp{\Gamma_2(f) - \rho \Gamma(f)}.
    \end{align*}

\begin{remark}
Without the second order $\Gamma(\Gamma(f)) \le 4\Gamma(f) \cdot (\Gamma_2(f) - \rho \Gamma(f))$ inequality: \url{https://arxiv.org/pdf/2403.00969}.
\end{remark}
\section{Gaussian isoperimetric inequality on discrete hypercube}

\begin{lemma}[Bobkov inequality on discrete hypercube] 
For any function $f: \{-1,+1\}^n \to [0,1]$, it holds
\begin{align*}
    I\tp{\int f \dif \mu_n} \le \int \sqrt{I(f)^2 + \abs{\nabla f}^2} \dif \mu_n,
\end{align*}
where $\mu_n$ is the uniform distribution over Boolean hypercube.
\end{lemma}

\begin{proof}[Proof sketch]
    We assume the inequality holds when $n=1$:
    \begin{align*}
        I\tp{\frac{a+b}{2}} \le \frac{1}{2} \tp{\sqrt{I(a)^2 + \tp{\frac{b-a}{2}}^2} + \sqrt{I(b)^2 + \tp{\frac{b-a}{2}}^2}}.
    \end{align*}
    We prove by induction hypothesis:
    Note that
    \begin{align*}  
        \int \sqrt{I(f)^2 + \abs{\nabla f}^2} \dif \mu_{n+1} &= \frac{1}{2} \int \sqrt{I(f_0)^2 + \abs{\nabla f_0}^2 + (f_1-f_0)^2/4} + \sqrt{I(f_1)^2 + \abs{\nabla f_1}^2 + (f_1-f_0)^2/4} \dif \mu_n \\
        &\ge \frac{1}{2} \tp{\sqrt{I(a)^2 + \tp{\frac{b-a}{2}}^2} + \sqrt{I(b)^2 + \tp{\frac{b-a}{2}}^2}}\\
        &\ge I(f),
    \end{align*}
    where $a = \int f_0 \dif \mu_n$, $b = \int f_1 \dif \mu_n$ and we use the inequality \begin{align*}
    \E[\mu_n]{\sqrt{a^2+b^2+c^2}} \ge \sqrt{\tp{\E[\mu_n]{\sqrt{a^2+b^2}}}^2 + \tp{\E[\mu_n]{c}}^2}
    \end{align*}
\end{proof}

\begin{remark}
    It seems that $(\partial_t - L) F_t \le 0$ fails to hold.
\end{remark}

\begin{remark}
For biased product distribution $\mu$ and function $h: \{-1,+1\}^n \to \{0,1\}$, we have
\begin{align*}
    I(\E[\mu]{h}) \le \sqrt{2} \E[\nu]{\sqrt{\sum_{i} \Var[i]{h}}}.
\end{align*}
\end{remark}

\section{Eldan-Gross inequality}
Eldan-Gross type inequality relies on the following ingredients:
\begin{lemma}
    For any $f:\{-1,+1\}^n \to \{-1,+1\}$:
    \begin{align*}
        I\tp{\frac{1+P_t f}{2}} \le \sqrt{\pi(1-\e^{-2t})} P_t \abs{\nabla f}.
    \end{align*}
\end{lemma}

\begin{lemma}[Hypercontractivity]
    \begin{align*}
    \Var[]{P_t f} \le \tp{\Var[]{f}}^{1-\Theta(t)} \cdot \tp{\sum_i \mathrm{Inf}_i(f)^2}^{\Theta(t)},
    \end{align*}
    where $\Theta(t) = \frac{1-\exp\tp{-2t}}{1+\exp\tp{-2t}}$ and $\mathrm{Inf}_i (f) = \E[]{\abs{\partial_i f}}$.
\end{lemma}

\begin{theorem}[Gross-Eldan]
    For any $f: \{-1,+1\}^n \to \{-1,+1\}$,
    \begin{align*}
        \int \abs{\nabla f} \dif \mu_n \gtrsim \Var[\mu_n]{f} \cdot \sqrt{\log \tp{1+ \frac{\e }{\sum_j \mathrm{Inf}_j(f)^2}}}.
    \end{align*}
\end{theorem}

\begin{proof}
    By Lemma 3.1 and inequality $I(x) \gtrsim x(1-x) \sqrt{\log \frac{1}{x(1-x)}}$, it holds
    \begin{align*}
        \E[]{\abs{\nabla{f}}} \gtrsim \Var[]{f} \cdot \sqrt{\log \frac{\e}{\Var[]{f}}}.
    \end{align*}
    Similarly, we have
    \begin{align*}
        \E[]{\abs{\nabla f}} \gtrsim \Var[]{f}.
    \end{align*}
    Therefore, we may assume without loss of generality that $\Var[]{f} \ge 100 W(f) := 100 \sum_{j} \mathrm{Inf}_j(f)^2$, $\Var[]{f} \ge \sqrt{W(f)}$ and $W(f) \ge 1/100$.

    By Lemma 3.1 and Lemma 3.2, it holds
    \begin{align*}
        \E[]{\abs{\nabla f}} \gtrsim \frac{\Var[]{f} - \Var[]{P_t f}}{\sqrt{1- \e^{-2t}}} \gtrsim \frac{\Var[]{f} - (\Var[]{f})^{1-\Theta(t)} W(f)^{\Theta(t)}}{\sqrt{1-\e^{-2t}}}.
    \end{align*}
    Choose an appropriate $t$ to conclude the proof.
\end{proof}

\section{Bobkov inequality on hypercube via semigroup}

We prove the following (relaxed) version of the Bobkov inequality
\begin{align*}
    I\tp{\int f \dif \mu_n} \le \int \sqrt{I(f)^2 + 2\abs{\nabla f}^2} \dif \mu_n.
\end{align*}
\begin{proof}
    Let $L f = \frac{1}{2} \sum_{i=1}^n \tp{f(x^{\oplus i}) - f(x)}$ and $f_t = \e^{Lt} f$.
    It suffices to show that $(\partial_t - L) F_t \le 0$ for all $x \in \{-1,+1\}^n$, where
    \begin{align*}
        F_t = \sqrt{I(f_t)^2 + 2 \abs{\nabla f_t}^2} = \abs{V_t},
    \end{align*}
    where $V_t = (I(f_t),\sqrt{2} \cdot \partial_1 f_t,\ldots,\sqrt{2} \cdot \partial_n f_t)$.
    Note that
    \begin{align*}
        L F_t(x) &= \frac{1}{2}\sum_{i=1}^n \tp{|V_t|(x^{\oplus i}) - |V_t|(x)}\\
        &= \frac{1}{2} \sum_{i=1}^n \tp{\abs{V_t}(x^{\oplus i}) - \frac{\abs{V_t}^2(x)}{\abs{V_t}(x)}}\\
        &= \frac{1}{2} \sum_{i=1}^n \tp{\abs{V_t}(x^{\oplus i}) + \frac{\inner{V_t(x)}{L V_t(x)}}{n \abs{V_t}(x)} - \frac{\inner{V_t(x)}{V_t(x^{\oplus i})}}{\abs{V_t}(x)}}\\
        &= \frac{\inner{V_t}{LV_t}}{\abs{V_t}}(x) + \frac{1}{2} \sum_{i=1}^n \underbrace{\tp{\abs{V_t}(x^{\oplus i}) - \frac{\inner{V_t(x)}{V_t(x^{\oplus i})} }{\abs{V_t}(x)}}}_{:=D_i}.
    \end{align*}
    Furthermore,
\begin{align*}
\partial_t F_t(x) = \partial_t \sqrt{\sum_{i=0}^n V_{t,i}^2} = \frac{\inner{V_{t}}{\partial_t V_t}}{\abs{V_t}}(x). 
\end{align*}
Hence, we have
\begin{align*}
    (\partial_t - L) F_t(x) = \frac{\inner{V_t}{(\partial_t - L)V_t}}{\abs{V_t}} - \frac{1}{2} \sum_{i=1}^n D_i.
\end{align*}
By a straightforward calculation,
\begin{align*}
    \partial_t \partial_i f_t = \partial_i \partial_t f_t = \partial_i L f_t = L \partial_i f_t.
\end{align*}
Therefore, it holds
\begin{align*}
    (\partial_t - L) F_t(x) = \frac{I(f_t) \cdot (\partial_t - L)I(f_t)}{\abs{V_t}} - \frac{1}{2} \sum_{i=1}^n D_i.
\end{align*}
We bound these two terms respectively.

Note that $V_{t,i}(x) = - V_t(x^{\oplus i})$. 
Let $\widetilde{V_t}(x)$ be $V_t(x)$ except $\widetilde{V_{t,i}}(x) = -V_{t,i}(x)$
Therefore,
\begin{align*}
    D_i = \frac{2V_{t,i}(x)^2}{\abs{V_t}(x)} + \abs{V_t}(x^{\oplus i}) - \frac{\inner{\widetilde{V}_{t}(x)}{V_{t}(x^{\oplus i})}}{\abs{V_t}(x)} \ge \frac{2 V_{t,i}(x)^2}{\abs{V_t}(x)}.
\end{align*}
Therefore, $\frac{1}{2}\sum_{i=1}^n D_i \ge \abs{V_t}$. 
Finally, note that
\begin{align*}
    (\partial_t - L) I(f_t) = I'(f_t(x)) L f_t(x) - L I(f_t) = \frac{1}{2} \sum_{i=1}^n I'(f_t(x)) \tp{f_t(x^{\oplus i}) - f_t(x)} - (I(f_t(x^{\oplus i})) - I(f_t(x))).
\end{align*}
Let $x_i = f_t(x^{\oplus i})$ and $x = f_t(x)$. It holds
\begin{align*}
    (\partial_t - L) I(f_t) = -\frac{1}{2} \sum_{i=1}^n \tp{I(x_i) - I(x)} - I'(x)(x_i-x).
\end{align*}
By inequality
\begin{align*}
    -I(a)(I(b)-I(a) - I'(a)(b-a)) \le (b-a)^2,
\end{align*}
It holds
\begin{align*}
    I(f_t) \cdot (\partial_t - L) I(f_t) \le 2 \abs{\nabla f_t}^2 \le \abs{V_t}^2.
\end{align*}
\end{proof}

\section{Bobkov inequality on slices via semigroup}
Let $\mu_{n,k}$ be the uniform distribution over $\{x \in \{-1,+1\}^n: \sum_i x_i = k\}$.
\begin{align*}
    I\tp{\int f \dif \mu_{n,k}} \le \int \sqrt{I(f(x))^2 + \sum_{\text{feasible swap } a,b} (f(x^{-a+b}) - f(x))^2} \dif \mu_{n,k}(z).
\end{align*}

{\color{red} TODO:
\begin{enumerate}
    \item Can we establish Bobkov inequality in more general settings (say \url{https://www.cs.cmu.edu/~odonnell/papers/kkl-kk.pdf});
    \item Implications of Bobkov inequality for slices.
\end{enumerate}
}
\end{document}

